\scanpage{34}
où
\[\unsure{\varpi_0|\bar\omega_0}\in\Gamma(\Omega^{-1}\otimes\Omega^1_{X/S})=\operatorname{Hom}(\Omega,\Omega^1_{X/S})\]
est indépendant de $i$, et est le $\unsure{\varpi_0|\bar\omega_0}$ de tout à l'heure pour un splitting. On dit que $C$ est connexion \uline{régulière} si la section $\unsure{\varpi_0|\bar\omega_0}$ de $\Omega^{-1}\otimes\Omega^1_{X/S}$ est \uline{régulière}.

\deleted{(\unsure{terme que nous utiliserons par la suite sans autres})}
\inserted{très régulière si} \deleted{\unsure{l'on a}}
\[\unsure{\varpi_0|\bar\omega_0}:\Omega\xrightarrow{\sim}\Omega^1_{X/S}\quad!\quad)\]
\uline{Prop} \inserted{(car. nulle)} Supposons
\[\mathcal P=\widehat S=\mathcal O_{\widehat P},\quad\text{où }(P,\sigma)\]
associé à une suite ex. de $0\longrightarrow\Omega\longrightarrow E\longrightarrow\mathcal O_X\longrightarrow0$.

Soit \deleted{\unsure{Supposons}} $\operatorname{Conn}(P)\deleted{\unsure{^{*}}}$ l'ens des conn. \deleted{\unsure{régulières}} $c$ de $P$, \deleted{telles que $\hat c$} $\operatorname{Conn}(\widehat S)$ \inserted{l'ens des} conn. de $\widehat S$, d'où $c\longmapsto\hat c$
\[\operatorname{Conn}(P)\longrightarrow\operatorname{Conn}(\widehat S)\]
et soit $\operatorname{Conn}^{*}(\widehat S)\subset\operatorname{Conn}(\widehat S)$ formé des conn. \uline{régulières} \inserted{(\unsure{ou très régt|un peu fort})} de $\widehat S$, $\operatorname{Conn}^{*}(P)$ son image inverse dans $\operatorname{Conn}(P)$. Soit $G$ ($G_0$) le gpe des auto de $\widehat S$ augmenté qui sont l'identité mod $J^3$ ($J^4$).

\inserted{Alors,} L'application
\diagram{\begin{tikzcd}
\operatorname{Conn}^{*}(P)\times G_0\arrow[r] & \operatorname{Conn}^{*}(\widehat S)\\
(c,g)\arrow[r,mapsto] & g.\hat c
\end{tikzcd}}
est \unsure{bijective|injective} (\unsure{resp.|unq.} bijective)
\margin{De plus, \unsure{l'action de $G_0$ dans $\operatorname{Conn}^{*}(\widehat S)$} \deleted{\unsure{de}\ \illegible{2}} \unsure{est libre}.}
\figure{\begin{tikzpicture}[x=1cm,y=1cm]
\draw (0.8,2.6) .. controls (.65,1.9) and (.65,1.1) .. (.9,.95) .. controls (1.5,.6) and (1.6,1.3) .. (.9,.95) .. controls (.65,.75) and (.65,.35) .. (.65,0);
\draw[->] (-1.2,-.25)--(-.8,-.25);
\end{tikzpicture}}
\uline{Cor} \inserted{Soit $c\in\operatorname{Conn}^{*}(P)$,}
L'ens $G(c)\subset G$ des $g\in G$ tels que
\[g\hat c\in\operatorname{Im}\operatorname{Conn}^{*}(P)\]
est tel que
\[
G(c)\xrightarrow{\sim}G/G_0\quad\begin{matrix}\simeq\\\Gamma(X,\Omega^{\otimes2})\end{matrix}
\]
% archive foliation: 33
